\chapter{Experimental Methodology}
\label{ch:method}

\section{Models, Hyperparameters, and Evaluation Protocol}

Four supervised classifiers span the spectrum of complexity: logistic
regression (linear baseline), random forest (bagging ensemble), XGBoost
and LightGBM (gradient boosting). Two unsupervised baselines (Isolation
Forest, LOF) are run on the binary day task to compare against
label-free methods. Deep models are excluded because tree ensembles match
or exceed them on tabular flow features~\cite{Grinsztajn2022} and the
additional complexity would crowd out the adversarial and operational
analyses. Mirsky et al.~\cite{Mirsky2018} demonstrate that an
unsupervised autoencoder ensemble (Kitsune) can match supervised
baselines under streaming conditions; integrating their approach is left
to future work.

Hyperparameters were chosen from documented library defaults plus minor
coarse tuning on the stratified validation split. The full configuration
is in Appendix~A; the most relevant values are \texttt{n\_estimators=300}
for random forest, \texttt{max\_depth=8} and
\texttt{learning\_rate=0.1} for the gradient-boosting libraries, and a
\texttt{saga} solver with maximum 2000 iterations and balanced class
weights for logistic regression.

All metrics are reported on the held-out test split, never on validation.
Validation is used solely for threshold tuning via Youden's~J and early
stopping. Stratified results use the held-out stratified test partition;
day results use Friday; LODO results aggregate across the four
attack-bearing folds (Tuesday, Wednesday, Thursday, Friday; the Monday
fold is benign-only and contributes only an FPR measurement). UNSW
results use the official test split.

All headline binary day-split numbers carry 95\% percentile bootstrap
confidence intervals ($B = 1000$ resamples, seed 42). Pairwise model
comparisons use McNemar's test with Yates continuity correction;
every pairwise comparison is significant at $p < 10^{-4}$. Reproducibility:
every result comes from a single \texttt{make full} run that takes
45--60 minutes on a 2024 MacBook Pro M-series. The random seed is fixed
(\texttt{random\_state = 42}). Multi-threaded tree ensembles introduce
approximately 1\% F1 drift between runs because of non-deterministic
thread scheduling at leaf-split time.

\section{Design Decisions}
\label{sec:design}

Many of the most consequential choices in this project are not modelling
choices but framing choices. Table~\ref{tab:design-decisions} records the
choices made alongside the alternatives considered.

\begin{table}[H]
    \centering
    \caption{Design decisions: alternatives considered and rationale.}
    \label{tab:design-decisions}
    \footnotesize
    \begin{tabular}{p{2.6cm}p{4.0cm}p{2.6cm}p{4.4cm}}
        \toprule
        \textbf{Decision} & \textbf{Alternatives considered} & \textbf{Chosen} & \textbf{Rationale} \\
        \midrule
        Dataset & CICIDS2017, UNSW-NB15, CSE-CIC-IDS2018, NSL-KDD, KDDCup99 & CICIDS2017 + UNSW-NB15 & Modern benchmark plus independent testbed for cross-validation. \\
        \addlinespace
        Split protocol & Stratified random, K-fold CV, day-based, LODO & All three: stratified, day-based, LODO & Reporting all three exposes the gap between published-style and realistic protocols. \\
        \addlinespace
        Task framing & Multi-class, binary, hierarchical, anomaly only & Both multi-class and binary & Multi-class is standard; binary is deployable; comparing them is the headline finding. \\
        \addlinespace
        Models & LogReg, RF, XGBoost, LightGBM, SVM, MLP, LSTM, transformer & LogReg + RF + XGBoost + LightGBM & Tree ensembles dominate tabular benchmarks; linear baseline anchors comparison; deep models out of scope. \\
        \addlinespace
        Threshold method & 0.5, F1-optimal, Youden's J, FPR-constrained & Youden's J + FPR-constrained & Prevalence-invariant; matches real SOC alert-budget constraint. \\
        \addlinespace
        Metric & Accuracy, F1, macro F1, ROC-AUC, recall@FPR & Macro F1 + recall@FPR & Macro F1 weights classes equally; recall@FPR is deployment-relevant. \\
        \addlinespace
        Adversarial attack & FGSM, PGD, C\&W, greedy coordinate-ascent & Greedy coordinate-ascent (score-query black-box) & Tree ensembles non-differentiable; query-only attacks realistic for tabular IDS. \\
        \addlinespace
        Defence & Madry adv-training, randomisation, certified, ensembles & Naïve adversarial training & Simplest baseline; reporting that it fails is itself informative. \\
        \addlinespace
        Triage layer & Static MITRE, learned mapping & Static MITRE ATT\&CK & Realistic for BSc scope; learned mapping needs telemetry beyond available. \\
        \bottomrule
    \end{tabular}
\end{table}
